% ============================================================================
% TEMPLATE DE CURRÍCULO — INSTRUÇÕES PARA A IA E PARA REVISÃO HUMANA
% ============================================================================
% FINALIDADE: adaptar informações reais à vaga, em português do Brasil.
% Comentários orientam o preenchimento; não aparecem no PDF nem executam regras.
% Os campos entre colchetes são estrutura de rascunho, nunca fatos sobre a pessoa.
%
% IA_FATOS — Use somente o currículo base confirmado e informações adicionais
% explicitamente fornecidas pela pessoa. A vaga indica prioridades, não fatos
% sobre o candidato.
%
% IA_AUSENCIAS — Algo não mencionado é "não informado", não "inexistente".
% Peça esclarecimento quando isso mudar um fato. Não deduza vínculo atual, datas,
% local ou domínio técnico. Remova campos opcionais sem suporte. Se for necessário
% deixar estrutura para preenchimento, identifique-a como PENDENTE e solicite os
% dados à pessoa. Não apresente um exemplo como experiência ou projeto realizado.
%
% IA_RASCUNHO — Ao devolver pendências, liste os dados que faltam e peça seu
% preenchimento. Um documento com PENDENTE ou colchetes não preenchidos é um
% rascunho, não uma versão final para candidatura. Antes do uso, a pessoa deve
% completar os campos com informações reais ou remover a entrada/seção.
% Não mantenha Lorem ipsum, exemplos de contatos ou links fictícios na versão final.
%
% IA_VAGA — Priorize evidências relevantes. Melhore clareza e organização sem
% modificar o sentido. Descreva ação, contexto e resultado somente se confirmados.
% Sem resultado informado, descreva a contribuição sem criar impacto ou números.
% Requisitos da vaga sem evidência ficam em uma lista de lacunas separada; não
% devem aparecer como habilidades da pessoa. Não prometa aprovação ou nota ATS.
%
% IA_ORDEM — Ordem padrão: cabeçalho, Experiência, Habilidades, Educação.
% Para primeira oportunidade ou estágio, mova Educação para após o cabeçalho.
% Destaque Projetos quando forem a melhor evidência prática disponível.
% Sem experiência profissional informada, omita Experiência e use Educação,
% Projetos e Habilidades. Não transforme atividade acadêmica em emprego.
% Não há número fixo de vínculos para "pouca experiência": considere contexto,
% duração e relevância confirmados; na dúvida, pergunte à pessoa.
%
% IA_LATEX — Preserve classe, pacotes, margens, tipografia e a macro \cventry.
% Altere campos, quantidade e ordem dos blocos conforme estas regras.
% Para ativar Projetos, remova o primeiro sinal de comentário de TODAS as linhas
% do bloco, inclusive os ambientes e suas terminações. Não deixe ambientes abertos.
% Escape caracteres especiais dos dados: &, porcentagem, $, #, _, {, }, ~, ^
% e barra invertida. Não escape os próprios comandos do template.
% Não acrescente comandos para executar programas ou ler arquivos externos.
%
% IA_REVISAO — Confira fatos, evidências, datas e números. A revisão humana é
% obrigatória antes do uso, inclusive quando não houver campos pendentes.
% ============================================================================

% --- TIPO E FORMATO DO DOCUMENTO ---
\documentclass[a4paper,10pt]{article}

% --- PACOTES: BASE / IDIOMA ---
\usepackage[utf8]{inputenc}
\usepackage[T1]{fontenc}
\usepackage[brazil]{babel}

% --- PACOTES: LAYOUT / TIPOGRAFIA ---
\usepackage{geometry}
\usepackage{parskip}
\usepackage{microtype}
\usepackage{enumitem}
\usepackage{titlesec}
\usepackage{array}
\usepackage{tabularx}

% --- PACOTES: LINKS / BOOKMARKS ---
\usepackage{hyperref}
\usepackage{bookmark}
\usepackage{xurl}

% --- CONFIGURAÇÃO: MARGENS / ESPAÇAMENTO ---
\geometry{top=1.5cm, bottom=1.5cm, left=1.5cm, right=1.5cm}
\setcounter{secnumdepth}{0}
\setlist[itemize]{
    leftmargin=0.75em,
    itemsep=0.2em,
    topsep=0.25em,
    parsep=0em,
    partopsep=0em
}
\pagestyle{empty}

% --- CONFIGURAÇÃO: METADADOS / LINKS ---
\hypersetup{
    pdftitle={Currículo},
    colorlinks=true,
    linkcolor=black,
    urlcolor=black,
    citecolor=black,
    bookmarksdepth=2
}

% --- CONFIGURAÇÃO: TÍTULOS ---
\titleformat{\section}
{\Large\bfseries}
{}
{0em}
{}
[\titlerule\vspace{0.5ex}]
\titleformat{\subsection}
{\normalsize\bfseries}
{}
{0em}
{}
\titlespacing*{\subsection}{0pt}{0.35em}{0.15em}

% --- MACRO: ENTRADA COM BOOKMARK ---
% Argumentos: {organização ou projeto}{local ou natureza}{cargo, curso ou papel}{período}.
% IA: campos opcionais não informados recebem {}. Não crie dados para ocupar
% células. Na versão final, não mantenha colchetes nem a palavra PENDENTE.
\newcounter{cventry}
\newcommand{\cventry}[4]{%
    \refstepcounter{cventry}%
    \phantomsection%
    \pdfbookmark[2]{#1}{cventry-\thecventry}%
    \noindent\begin{tabularx}{\textwidth}{@{}>{\raggedright\arraybackslash}X >{\raggedleft\arraybackslash}X@{}}
    \textbf{#1} & #2 \\
    \textit{#3} & \textit{#4} \\
    \end{tabularx}
}

% --- INÍCIO DO DOCUMENTO ---
\begin{document}

% --- CABEÇALHO ---
% IA: use apenas contatos fornecidos. Não construa URLs a partir do nome.
% Remova contato opcional ausente e seu separador, sem deixar marcadores isolados.
% Os endereços abaixo exemplificam a sintaxe; não são contatos reais.
% Sem nome informado, mantenha o campo no rascunho e solicite-o à pessoa.
\begin{center}
    {\LARGE \textbf{[Nome completo]}}
    \\ [0.1cm]
    [Cidade, UF]
    {\textbullet}
    [Telefone]
    {\textbullet}
    \href{mailto:seu.email@exemplo.com}{[E-mail]}
    {\textbullet}
    \href{https://linkedin.com/in/seu-perfil}{[LinkedIn]}
    {\textbullet}
    \href{https://github.com/seu-perfil}{[Portfólio ou GitHub]}
\end{center}

% IA_PRIMEIRA_OPORTUNIDADE: mova o bloco inteiro de Educação para este ponto.
% Mova também seus comentários e remova a cópia no fim, sem duplicar a seção.

% --- EXPERIÊNCIA ---
% IA: use ordem cronológica inversa e somente vínculos/atividades confirmados.
% Estágio, monitoria, pesquisa ou voluntariado devem ter sua natureza identificada.
% Os quatro blocos são exemplos de estrutura, não uma quantidade a preencher.
% Mantenha só entradas com evidência e repita a estrutura quando necessário.
% Use "Atual" apenas se a continuidade do vínculo tiver sido confirmada.
% Sem experiência informada: omita a seção inteira e aplique IA_PROJETOS.
% Com pouca experiência relevante: preserve vínculos reais e destaque Projetos.
\section{Experiência}
    \cventry{[Empresa ou instituição]}{[Cidade, UF]}{[Cargo ou atividade]}{[Início] -- [Fim]}
        \begin{itemize}
            \item \textit{[Descreva uma atividade realizada e sua contribuição concreta.]}
            \item \textit{[Descreva uma entrega ou solução, com ferramentas comprovadas.]}
            \item \textit{[Inclua um resultado somente se houver evidência no currículo base.]}
        \end{itemize}

    \cventry{[Empresa ou instituição]}{[Cidade, UF]}{[Cargo ou atividade]}{[Início] -- [Fim]}
        \begin{itemize}
            \item \textit{[Descreva a contribuição mais relevante para a vaga.]}
            \item \textit{[Acrescente outra atividade real, sem repetir informações.]}
        \end{itemize}

    \cventry{[Empresa ou instituição]}{[Cidade, UF]}{[Cargo ou atividade]}{[Início] -- [Fim]}
        \begin{itemize}
            \item \textit{[Informe uma responsabilidade confirmada e seu contexto.]}
            \item \textit{[Informe uma entrega ou aprendizagem prática comprovada.]}
        \end{itemize}

    \cventry{[Empresa ou instituição]}{[Cidade, UF]}{[Cargo ou atividade]}{[Início] -- [Fim]}
        \begin{itemize}
            \item \textit{[Descreva uma atividade real sem criar cargo ou resultado.]}
            \item \textit{[Mantenha apenas detalhes relevantes e sustentados pela fonte.]}
        \end{itemize}

% --- PROJETOS: BLOCO OPCIONAL, INICIALMENTE COMENTADO ---
% IA_PROJETOS_ATIVAR — Para primeira oportunidade, estágio ou pouca experiência
% profissional relevante, descomente o bloco e destaque projetos pessoais,
% acadêmicos ou profissionais efetivamente informados. Inclua TCC, extensão ou
% trabalho de disciplina apenas se mencionados. Não crie projetos a partir da vaga.
% Mesmo com experiência, mantenha projetos reais se acrescentarem evidência útil.
%
% IA_PROJETOS_AUSENTES — Nesse perfil, se nenhum projeto estiver informado,
% descomente a seção e mantenha UMA entrada de estrutura, com o aviso PENDENTE.
% Peça nome, objetivo, contribuição individual, tecnologias, período e link,
% quando existirem. Não preencha campos com dados presumidos.
% Pergunte se a pessoa tem algum projeto que possa inserir. Se ela confirmar que
% ainda não tem, explique que a estrutura é um guia para preenchimento futuro,
% e não um projeto já realizado. A pessoa deve completar ou remover a pendência
% na revisão antes de usar o currículo em uma candidatura.
%
% IA_PROJETOS_PREENCHIDOS — Havendo projetos confirmados, remova o aviso PENDENTE,
% preencha apenas dados reais e duplique a entrada se necessário. Campos opcionais
% ausentes ficam {}. Não mantenha uma entrada vazia junto das entradas preenchidas.
%
% \section{Projetos}
%     \textit{PENDENTE: informe um projeto real para preencher a estrutura abaixo.
%     Se ainda não houver um, esta entrada é somente um guia de preenchimento.}
%     \cventry{[Nome do projeto]}{[Natureza do projeto]}{[Sua participação]}{[Período]}
%         \begin{itemize}
%             \item \textit{[Informe o problema ou objetivo e a contribuição individual.]}
%             \item \textit{[Informe tecnologias e entregas efetivamente realizadas.]}
%             \item \textit{[Inclua resultado ou link apenas se tiver sido informado.]}
%         \end{itemize}

% --- HABILIDADES ---
% IA: priorize competências comprovadas e pertinentes à vaga, sem copiar todos
% os requisitos. Ter cursado um assunto não comprova domínio avançado.
% Preserve o nível de idioma informado; sem nível, não presuma fluência.
% Sem idiomas informados, remova essa linha. Sem habilidades sustentadas, omita
% a seção e solicite os dados na revisão.
\section{Habilidades}
    \begin{itemize}
        \item \textbf{Técnicas:} [Competências, ferramentas e tecnologias confirmadas]
        \item \textbf{Idiomas:} [Idiomas e níveis de proficiência informados]
    \end{itemize}

% --- EDUCAÇÃO ---
% IA_EDUCACAO — Para primeira vaga ou estágio, mova este bloco para após o cabeçalho.
% Nos demais casos, preserve a posição padrão, salvo justificativa de relevância.
% Liste a formação mais recente primeiro e repita apenas para cursos informados.
% Preserve instituição, curso e situação. "Conclusão prevista" só pode ser usada
% quando a previsão for confirmada. Não trate formação em andamento como concluída.
% Local e data são opcionais: use {} quando faltarem. Sem formação informada,
% solicite instituição, curso e situação; não invente uma formação para preencher.
\section{Educação}
    \cventry{[Instituição de ensino]}{[Cidade, UF]}{[Curso ou formação]}{[Situação e período]}

% IA_CHECAGEM_FINAL — Procure PENDENTE, colchetes não preenchidos, links de exemplo,
% seções vazias e informações sem fonte. Havendo pendências, devolva como rascunho
% e peça os dados à pessoa; nunca apresente esses campos como informações finais.
\end{document}
